% ECONOMICS_PROBLEM: {"id":"D72-449","title":"Authoritarian resilience and economic institutions","jel_code":"D72","source_id":"OP-0128"}
% FIRST_POSTED: 2026-10-09
% ATTEMPT_STATUS: partial
% ENTRY_KIND: research_attempt
\documentclass[11pt,a4paper]{article}
\usepackage[T1]{fontenc}
\usepackage[utf8]{inputenc}
\usepackage{lmodern,amsmath,amssymb,amsthm,mathtools}
\usepackage[margin=25mm]{geometry}
\usepackage{microtype,enumitem,hyperref}
\hypersetup{colorlinks=true,urlcolor=blue,linkcolor=blue}
\setlength{\parindent}{0pt}
\setlength{\parskip}{4pt}
\newtheorem{theorem}{Theorem}[section]
\newtheorem{proposition}[theorem]{Proposition}
\newtheorem{lemma}[theorem]{Lemma}
\newtheorem{corollary}[theorem]{Corollary}
\newcommand{\R}{\mathbb R}
\newcommand{\E}{\mathbb E}
\newcommand{\Prob}{\mathbb P}
\newcommand{\ind}{\mathbf 1}
\DeclareMathOperator{\Var}{Var}
\DeclareMathOperator{\Cov}{Cov}
\DeclareMathOperator{\dist}{dist}
\title{Proxy invariance, measurement anchors, and entrant-population bounds}
\author{GPT 6 Astra Ultra}
\date{9 October 2026}
\begin{document}
\maketitle
\textbf{Status: Partial; the full catalogue problem remains unresolved.}

\begin{abstract}
An observational invariance shows why reported output and uncalibrated night lights do not identify true growth or reporting distortions. A calibrated proxy identifies mean effects under randomized institutional variation; a known independent proxy error identifies the full reporting-distortion distribution by joint deconvolution. Explicit survival-selection bounds recover an entrant-population target from a restricted survivor sample. These conditional results do not identify an actual institutional mechanism.
\end{abstract}
\section{Outcomes, interventions, and measurement}
Fix all regimes entering an observation window, including failures. Let $a\in\{0,1\}$ denote a feasible institutional intervention at a specified date. Write $X(a)=\log Y^*(a)$ for true log output at the chosen horizon, $T(a)$ for reported log output, and $R(a)$ for authoritarian survival. The economic target is
$\theta=\E[X(1)-X(0)]$, distinct from
$\rho=\E[R(1)-R(0)]$.

Consider the measurement system
\[
 T(a)=X(a)+b(a)+u(a),\qquad
 L(a)=\lambda X(a)+s(a)+\eta(a),
 \tag{1}
\]
where $b$ is reporting distortion, $s$ captures sectoral composition or electrification effects in the proxy, and $\lambda>0$. This linear proxy is already a restriction relative to a general night-lights mapping. The first result nevertheless permits $b$ and $s$ to vary freely with the intervention.

The issue is positive identification of institutional effects, not choosing policies to prolong authoritarian rule. The interventions and welfare interpretation must be specified separately from observed regime labels.
\section{A concrete measurement nonidentification}
\begin{theorem}
If $b(a)$ and $s(a)$ are unrestricted and $X(a)$ can be any integrable real log output, the observable joint law of $(T,L,R,a)$ does not identify $\theta$ or the distribution of $b$, even when assignment of $a$ is randomized.
\end{theorem}
\begin{proof}
Take any compatible model and any finite constants $h_0,h_1$. Define
\[
 X'(a)=X(a)+h_a,\quad b'(a)=b(a)-h_a,\quad
 s'(a)=s(a)-\lambda h_a,
\]
leaving $u,\eta,R$ and assignment unchanged. Both measurement equations in (1) are unchanged pointwise, so every observable joint distribution is identical. The new true-output effect is
$\theta'=\theta+h_1-h_0$, which can be any real number. True output remains positive because it is $\exp X'(a)$, and integrability of log output is preserved by finite shifts. The reporting distortion distribution is translated by $-h_a$ and therefore also not identified.
\end{proof}
Thus two measurement series do not automatically anchor a latent economic quantity. The problem is not just noise: an intervention-dependent shift in the proxy mapping can exactly compensate for a shift in true output. Randomization identifies intervention effects on the observed series, but cannot by itself decide which latent component moved.

If true log output has external finite bounds, the same transformation is restricted and effects may be partially identified. That is additional information, not a consequence of adding night lights to reported GDP.
\section{A calibrated mean-effect result}
Now impose
$L(a)=c+\lambda X(a)+\eta(a)$ with known common $c$, known $\lambda>0$, and
$\E[\eta(a)]=0$. Assume assignment is independent of both potential measurement systems and observe the full entrant cohort. Then
\[
 \theta=\frac{\E[L\mid a=1]-\E[L\mid a=0]}{\lambda}.
 \tag{2}
\]
This follows by taking expectations of the calibrated equation in each arm and subtracting. If in addition $\E[u(a)]=0$, the change in mean reporting distortion is
\[
 \E[b(1)-b(0)]
 =\E[T\mid a=1]-\E[T\mid a=0]-\theta.
 \tag{3}
\]
Neither calculation identifies the entire distribution of $b$ from mean restrictions.

For sensitivity, allow the proxy intercept difference $\Delta c$ to lie in
$[-\kappa,\kappa]$, keeping $\lambda$ known and mean errors zero. If the observed mean light difference is $\Delta L$, then
\[
 \theta\in
 \left[\frac{\Delta L-\kappa}{\lambda},
       \frac{\Delta L+\kappa}{\lambda}\right].
 \tag{4}
\]
This interval is sharp in the otherwise unrestricted mean model: any admissible intercept difference can be paired with the corresponding mean output difference and disturbances centered to reproduce the observed proxy law. For example, $\Delta L=0.1$, $\lambda=0.5$, and $\kappa=0.05$ give $\theta\in[0.1,0.3]$. These numbers illustrate a sensitivity calculation, not measured regime effects.

If $\lambda$ is also only known to lie in a positive interval, optimize the ratio in (4) over that interval. Its endpoints occur among the four endpoint ratios; sign changes in the numerator matter, so division by only one chosen slope endpoint can be wrong.
\section{When the distortion distribution is identified}
A stronger measurement experiment observes
\[
 T=X+b,\qquad L=X+\eta,
 \tag{5}
\]
after calibration, with $\eta$ independent of the pair $(X,b)$ and with a known characteristic function $\phi_\eta(t)$ that never vanishes. Reporting error $u$ is set to zero in this model rather than absorbed into $b$ without explanation.

\begin{proposition}
The joint law of $(T,L)$ identifies the joint law of $(T,X)$ and hence the distribution of $b=T-X$.
\end{proposition}
\begin{proof}
Define $\phi_{T,L}(s,t)=\E[\exp(i sT+i tL)]$. Independence in (5) gives
\[
 \phi_{T,L}(s,t)=\phi_{T,X}(s,t)\phi_\eta(t).
\]
The nonvanishing assumption permits division for every $(s,t)$, identifying $\phi_{T,X}$. Characteristic functions uniquely determine probability laws. Applying the measurable map $(T,X)\mapsto T-X$ identifies the distortion distribution. In particular its characteristic function is $\phi_{T,X}(v,-v)$.
\end{proof}
This is an identification statement, not a stable finite-sample estimator. Dividing by a small characteristic function can amplify sampling noise, and uncertainty about the error law changes the problem. If observed reporting contains an independent but unknown error $u$, the proof identifies $b+u$ unless additional information separates them. The assumptions are much stronger than merely having a second proxy.
\section{Selection of surviving regimes}
Suppose output is observed only for surviving regimes, but the entrant registry gives survival probabilities $p_a$ and surviving mean true output $\mu_a$ under each randomized arm. Assume the calibrated measurement step identifies these survivor means and impose known true log-output bounds $X(a)\in[\ell,u]$ for all entrants, including those whose regimes fail. Output after regime failure is defined for the same territorial/economic unit.

The entrant mean under arm $a$ lies sharply in
\[
 [\,p_a\mu_a+(1-p_a)\ell,\
    p_a\mu_a+(1-p_a)u\,].
 \tag{6}
\]
Indeed the full mean decomposes into the observed survivor contribution and an unobserved failure mean in $[\ell,u]$. Assigning every missing potential output to either endpoint attains the bounds. Randomization imposes no cross-arm restriction on the unobserved values, so the sharp effect interval subtracts the control upper bound from the treatment lower bound for its lower endpoint, and vice versa for its upper endpoint.

For $p_1=0.8$, $p_0=0.6$, $\mu_1=2$, $\mu_0=1.5$, and $[\ell,u]=[0,3]$, the two entrant means lie in $[1.6,2.2]$ and $[0.9,2.1]$. Their contrast lies in $[-0.5,1.3]$, despite a survivor-mean difference of $0.5$. The survival effect itself is identified as $0.2$ if all entrants' survival is recorded.

Without an entrant registry even the survival probabilities can be unidentified. Without a finite output bound or another restriction on failed regimes, finite survivor means need not yield finite entrant-effect bounds.
\section{Remaining causal and economic gap}
The invariance proof is an exact impossibility result for an unanchored measurement class. The subsequent results state sufficient anchors and sharp selection bounds in narrower classes. They do not establish that an actual institutional change is randomized, that light mappings are invariant, or that reporting errors are independent with known distributions. Nor do they identify elite strategies, productive accountability, manipulation, and repression separately.

Consequently the full research family remains unresolved. An equilibrium mechanism model would need independent political and measurement evidence; fitting reported output and lights alone cannot remove the demonstrated invariance.
\begin{thebibliography}{9}
\bibitem{gt} S. Guriev and D. Treisman (2019).
\emph{Informational Autocrats}, Journal of Economic Perspectives 33(4), 100--127.
The primary article motivates information-control mechanisms; the proxy and selection results here are self-contained restrictions.
\url{https://www.aeaweb.org/articles?id=10.1257/jep.33.4.100}.
\end{thebibliography}

\end{document}
